\documentclass[aspectratio=169, 10pt]{beamer}

% --- Theme & Color Configuration ---
\usetheme{Madrid}
\usecolortheme{whale}

\definecolor{primaryDark}{RGB}{30, 41, 59}     % Slate 800
\definecolor{accentBlue}{RGB}{37, 99, 235}     % Blue 600
\definecolor{bgGray}{RGB}{248, 250, 252}       % Slate 50

\setbeamercolor{palette primary}{bg=primaryDark, fg=white}
\setbeamercolor{palette secondary}{bg=accentBlue, fg=white}
\setbeamercolor{palette tertiary}{bg=primaryDark, fg=white}
\setbeamercolor{titlelike}{bg=primaryDark, fg=white}
\setbeamercolor{block title}{bg=accentBlue, fg=white}
\setbeamercolor{block body}{bg=bgGray, fg=black}

\setbeamertemplate{navigation symbols}{}
\setbeamertemplate{footline}[frame number]

% --- Font & Language Setup ---
\usepackage{fontspec}
\usepackage[english, bidi=basic, provide=*]{babel}
\babelprovide[import, onchar=ids fonts]{portuguese}
\babelfont{rm}{Noto Sans}
\babelfont{sf}{Noto Sans}

\usepackage{amsmath, amssymb}
\usepackage{booktabs}
\usepackage{listings}

\title[Criptografia RSA]{Criptografia RSA}
\subtitle{Fundamentação Teórica e Implementação Prática}
\author{Segurança da Informação / Criptografia Aplicada}
\date{\today}

\begin{document}

% --- Slide 1: Title ---
\begin{frame}
    \titlepage
\end{frame}

% --- Slide 2: Visão Geral ---
\begin{frame}{Visão Geral do Conteúdo}
    \begin{columns}[T]
        \begin{column}{0.48\textwidth}
            \begin{block}{1. Conceito e História}
                \begin{itemize}
                    \item Criptografia simétrica vs. assimétrica
                    \item Origem histórica (Diffie-Hellman, RSA)
                \end{itemize}
            \end{block}
            \begin{block}{2. Matemática do RSA}
                \begin{itemize}
                    \item Geração de chaves passo a passo
                    \item Exemplo numérico e fórmulas
                \end{itemize}
            \end{block}
        \end{column}
        \begin{column}{0.48\textwidth}
            \begin{block}{3. Segurança e Ameaças}
                \begin{itemize}
                    \item Dificuldade de fatoração
                    \item Tamanho de chaves e Criptografia Quântica
                \end{itemize}
            \end{block}
            \begin{block}{4. Aplicação Prática}
                \begin{itemize}
                    \item Arquitetura do código em Python
                    \item Chat seguro via Sockets TCP
                \end{itemize}
            \end{block}
        \end{column}
    \end{columns}
\end{frame}

% --- Slide 3: Por que Chave Pública? ---
\begin{frame}{Criptografia Simétrica vs. Assimétrica}
    \begin{table}
        \centering
        \begin{tabular}{lp{5.5cm}p{5.5cm}}
            \toprule
            \textbf{Característica} & \textbf{Simétrica (ex.: AES)} & \textbf{Assimétrica / RSA} \\
            \midrule
            \textbf{Chaves} & Uma única chave secreta compartilhada & Par de chaves: Pública e Privada \\
            \textbf{Vantagem} & Alta velocidade de processamento & Elimina a troca secreta prévia de chaves \\
            \textbf{Desvantagem} & Desafio seguro na distribuição da chave & Desempenho computacional mais lento \\
            \textbf{Aplicação} & Cifragem de grandes volumes de dados & Troca de chaves de sessão e assinaturas \\
            \bottomrule
        \end{tabular}
    \end{table}
    \vspace{0.3cm}
    \centering
    \textbf{Princípio do RSA:} Cifra-se com a \textit{chave pública} do destinatário; apenas a \textit{chave privada} correspondente decifra.
\end{frame}

% --- Slide 4: Matemática do RSA ---
\begin{frame}{Geração de Chaves em 5 Passos}
    \begin{enumerate}
        \item \textbf{Seleção de Primos:} Escolher dois primos grandes e distintos, $p$ e $q$.
        \item \textbf{Módulo do RSA ($n$):} Calcular $n = p \cdot q$ (parte pública da chave).
        \item \textbf{Totiente de Euler ($\phi$):} Calcular $\phi(n) = (p - 1)(q - 1)$.
        \item \textbf{Expoente Público ($e$):} Escolher $e$ tal que $1 < e < \phi(n)$ e $\text{mdc}(e, \phi(n)) = 1$.
        \item \textbf{Expoente Privado ($d$):} Calcular $d \equiv e^{-1} \pmod{\phi(n)}$ (inverso modular).
    \end{enumerate}

    \begin{columns}[T]
        \begin{column}{0.48\textwidth}
            \begin{block}{Chave Pública}
                $(e, n)$ --- Pode ser divulgada abertamente.
            \end{block}
        \end{column}
        \begin{column}{0.48\textwidth}
            \begin{block}{Chave Privada}
                $(d, n)$ --- Deve ser mantida sob segredo absoluto.
            \end{block}
        \end{column}
    \end{columns}
\end{frame}

% --- Slide 5: Exemplo Numérico ---
\begin{frame}{Exemplo Numérico e Cifragem}
    \begin{columns}[T]
        \begin{column}{0.48\textwidth}
            \begin{block}{Parâmetros do Exemplo}
                \begin{itemize}
                    \item Primos: $p = 61, \quad q = 53$
                    \item Módulo: $n = 61 \times 53 = 3233$
                    \item Totiente: $\phi(n) = 60 \times 52 = 3120$
                    \item Exp. Público: $e = 17$
                    \item Exp. Privado: $d = 2753$
                \end{itemize}
            \end{block}
        \end{column}
        \begin{column}{0.48\textwidth}
            \begin{block}{Cifragem e Decifragem}
                Para uma mensagem $m = 65$:
                \begin{align*}
                    c &= m^e \pmod n = 65^{17} \pmod{3233} = 2790 \\[4pt]
                    m &= c^d \pmod n = 2790^{2753} \pmod{3233} = 65
                \end{align*}
            \end{block}
            \small
            \textbf{Validação:} $(m^e)^d \equiv m^{e \cdot d} \equiv m \pmod n$ graças ao Teorema de Euler.
        \end{column}
    \end{columns}
\end{frame}

% --- Slide 6: Segurança e Tamanho de Chaves ---
\begin{frame}{Segurança, Chaves e Ameaçadoras Futuras}
    \begin{columns}[T]
        \begin{column}{0.48\textwidth}
            \begin{block}{Recomendações de Chave}
                \begin{itemize}
                    \item \textbf{512 bits:} Inseguro (quebrado empiricamente).
                    \item \textbf{1024 bits:} Desaconselhado / Obsoleto.
                    \item \textbf{2048 bits:} Padrão mínimo recomendado atual.
                    \item \textbf{3072+ bits:} Seguro para longo prazo.
                \end{itemize}
            \end{block}
        \end{column}
        \begin{column}{0.48\textwidth}
            \begin{block}{Ameaças e Vulnerabilidades}
                \begin{itemize}
                    \item \textbf{Fatoração de $n$:} Base da segurança. Algoritmos modernos como GNFS exigem $O(e^c)$ operações.
                    \item \textbf{Criptografia Quântica:} O \textit{Algoritmo de Shor} fatora inteiros em tempo polinomial em computadores quânticos.
                    \item \textbf{Ataques Práticos:} Uso de RSA puro sem \textit{padding} (OAEP) permite exploração por maleabilidade.
                \end{itemize}
            \end{block}
        \end{column}
    \end{columns}
\end{frame}

% --- Slide 7: Implementação e Chat ---
\begin{frame}{Implementação do Chat Criptografado}
    \begin{columns}[T]
        \begin{column}{0.48\textwidth}
            \begin{block}{Componentes do Núcleo (`rsa\_core.py`)}
                \begin{itemize}
                    \item \textbf{Miller-Rabin:} Teste probabilístico de primalidade.
                    \item \textbf{Euclides Estendido:} Cálculo do inverso modular $d$.
                    \item \textbf{Square-and-Multiply:} Exponenciação modular rápida $a^b \pmod n$.
                \end{itemize}
            \end{block}
        \end{column}
        \begin{column}{0.48\textwidth}
            \begin{block}{Fluxo de Rede (Sockets TCP)}
                \begin{enumerate}
                    \item Conexão cliente-servidor estabelecida.
                    \item Handshake de chaves públicas.
                    \item Cifragem local das mensagens enviadas.
                    \item Decifragem em tempo real das mensagens recebidas.
                \end{enumerate}
            \end{block}
        \end{column}
    \end{columns}
\end{frame}

% --- Slide 8: Conclusão ---
\begin{frame}{Conclusão}
    \begin{itemize}
        \item \textbf{Fundamento Histórico:} O RSA resolveu o problema crítico da troca de chaves em redes abertas.
        \item \textbf{Dependência Matemática:} A segurança depende diretamente da complexidade computacional da fatoração de inteiros grandes.
        \item \textbf{Boas Práticas:} Em ambientes reais, o RSA deve ser combinado com esquemas de preenchimento seguro (OAEP) e certificados digitais (X.509/PKI).
        \item \textbf{Perspectiva de Futuro:} O avanço do processamento quântico motiva a transição contínua para algoritmos da Criptografia Pós-Quântica (PQC).
    \end{itemize}
\end{frame}

\end{document}